\mxunit{expertes-c11}{Suites de matrices et chaînes de Markov}
\mxsubtitle{Suites $U_{n+1}=AU_n+B$, graphes probabilistes, distributions invariantes}

\begin{objectives}
Étudier une suite de matrices colonnes définie par $U_{n+1}=AU_n+B$ ; modéliser une évolution aléatoire à deux ou trois états par un graphe probabiliste et une matrice de transition ; calculer la distribution après $n$ étapes ; déterminer une distribution invariante ; étudier la convergence dans le cas de deux états.
\end{objectives}
\begin{prerequisites}
Chapitres 9 et 10 ; probabilités conditionnelles et formule des probabilités totales ; suites arithmético-géométriques ; limite d'une suite géométrique.
\end{prerequisites}

\section{Automatismes}
\begin{fluency}[Sans calculatrice, cinq minutes]
\begin{enumerate}
\item On pose $u_{n+1}=\num{0,5}u_n+2$. Déterminer le réel $\ell$ tel que $\ell=\num{0,5}\ell+2$.
\item Donner la limite de $\num{0,8}^n$ quand $n$ tend vers $+\infty$.
\item Dans un arbre, $P(A)=\num{0,3}$, $P_A(B)=\num{0,6}$, $P_{\overline A}(B)=\num{0,2}$. Calculer $P(B)$.
\item Résoudre le système $\num{0,2}a=\num{0,4}b$ et $a+b=1$.
\item Calculer $\begin{pmatrix}1&0\end{pmatrix}\begin{pmatrix}\num{0,8}&\num{0,2}\\\num{0,4}&\num{0,6}\end{pmatrix}$.
\item Quelle est la somme des coefficients de chaque ligne de la matrice de la question 5 ?
\end{enumerate}
\end{fluency}
\begin{solution}[Automatismes]
\begin{enumerate}
\item $\ell=4$.
\item $0$.
\item $\num{0,3}\times\num{0,6}+\num{0,7}\times\num{0,2}=\num{0,32}$.
\item $a=2b$, donc $3b=1$ : $a=\frac23$ et $b=\frac13$.
\item $\begin{pmatrix}\num{0,8}&\num{0,2}\end{pmatrix}$.
\item $1$ pour chaque ligne.
\end{enumerate}
\end{solution}

\mxnewpage
\section{Activité d'introduction}
\begin{activity}[Beau temps, pluie \mxchip{40 min}]
Dans un modèle simplifié (données fictives), le temps d'une journée est soit beau ($S$), soit pluvieux ($P$), et ne dépend que du temps de la veille : après un beau jour, il fait beau avec la probabilité $\num{0,8}$ ; après un jour de pluie, il fait beau avec la probabilité $\num{0,4}$. Le jour $0$ est beau. On note $s_n$ et $p_n$ les probabilités que le jour $n$ soit beau ou pluvieux.
\begin{enumerate}
\item Représenter la situation par un schéma : deux sommets $S$ et $P$, une flèche pour chaque passage possible, portant sa probabilité.
\item À l'aide d'un arbre, calculer $s_1$, $p_1$, puis $s_2$, $p_2$.
\item Justifier que $s_{n+1}=\num{0,8}s_n+\num{0,4}p_n$ et $p_{n+1}=\num{0,2}s_n+\num{0,6}p_n$.
\item On pose $\pi_n=\begin{pmatrix}s_n&p_n\end{pmatrix}$ et $T=\begin{pmatrix}\num{0,8}&\num{0,2}\\\num{0,4}&\num{0,6}\end{pmatrix}$. Vérifier que $\pi_{n+1}=\pi_nT$. Calculer $\pi_{10}$ à la calculatrice. Recommencer en partant d'un jour $0$ pluvieux. Que constate-t-on ?
\item Déterminer la répartition $\pi=\begin{pmatrix}a&b\end{pmatrix}$, avec $a+b=1$, qui vérifie $\pi T=\pi$. La comparer aux observations de la question 4.
\end{enumerate}
\end{activity}
\begin{solution}[Activité]
\begin{enumerate}
\item Flèches $S\to S$ ($\num{0,8}$), $S\to P$ ($\num{0,2}$), $P\to S$ ($\num{0,4}$), $P\to P$ ($\num{0,6}$).
\item $s_1=\num{0,8}$, $p_1=\num{0,2}$ ; $s_2=\num{0,8}\times\num{0,8}+\num{0,2}\times\num{0,4}=\num{0,72}$ et $p_2=\num{0,28}$.
\item Formule des probabilités totales avec la partition « beau le jour $n$ », « pluie le jour $n$ ».
\item $\pi_nT=\begin{pmatrix}\num{0,8}s_n+\num{0,4}p_n&\num{0,2}s_n+\num{0,6}p_n\end{pmatrix}=\pi_{n+1}$. On trouve $\pi_{10}\simeq\begin{pmatrix}\num{0,6667}&\num{0,3333}\end{pmatrix}$, et $\pi_{10}\simeq\begin{pmatrix}\num{0,6666}&\num{0,3334}\end{pmatrix}$ en partant d'un jour pluvieux. La répartition semble se stabiliser vers $\begin{pmatrix}\frac23&\frac13\end{pmatrix}$, quel que soit le départ.
\item $\num{0,8}a+\num{0,4}b=a$ équivaut à $\num{0,2}a=\num{0,4}b$, soit $a=2b$ ; avec $a+b=1$ : $\pi=\begin{pmatrix}\frac23&\frac13\end{pmatrix}$, la limite observée.
\end{enumerate}
\end{solution}
\begin{teachernote}[Conduite de l'activité]
Les calculs numériques de la question 4 suggèrent la convergence ; la preuve, dans le cas de deux états, se fait avec une suite arithmético-géométrique (cours). Attirer l'attention sur l'écriture en ligne de $\pi_n$ et sur l'ordre du produit $\pi_nT$.
\end{teachernote}

\mxnewpage
\section{Notes de cours}
\subsection{Suites de matrices}
\begin{theorem}[Suite $U_{n+1}=AU_n+B$]
Soit $A$ une matrice carrée d'ordre $p$, $B$ une matrice colonne et $(U_n)$ une suite de matrices colonnes telle que $U_{n+1}=AU_n+B$ pour tout $n$. Si $I_p-A$ est inversible, la matrice $C=(I_p-A)^{-1}B$ vérifie $C=AC+B$, et pour tout entier naturel $n$ :
\[U_n=A^n\left(U_0-C\right)+C.\]
\end{theorem}
\begin{proof}
$(I_p-A)C=B$ s'écrit $C=AC+B$. En soustrayant : $U_{n+1}-C=A(U_n-C)$. La suite $V_n=U_n-C$ vérifie $V_{n+1}=AV_n$, donc $V_n=A^nV_0$ (récurrence), et $U_n=A^n(U_0-C)+C$.
\end{proof}
\begin{keypoint}[Comportement à l'infini]
Si tous les coefficients de $A^n$ tendent vers $0$, alors $U_n$ converge vers $C$ (coefficient par coefficient). Ce comportement s'observe à la calculatrice ; on l'admet lorsque l'énoncé le précise.
\end{keypoint}

\subsection{Chaînes de Markov}
\begin{definition}[Chaîne de Markov à deux ou trois états]
On observe un système à des instants $0,1,2,\dots$ ; à chaque instant, il se trouve dans un des états $1,\dots,k$ ($k=2$ ou $3$). On note $X_n$ l'état à l'instant $n$. On dit que $(X_n)$ est une \emph{chaîne de Markov} (homogène) si la probabilité de passer de l'état $i$ à l'état $j$ ne dépend ni de $n$ ni des états antérieurs : pour tous $i$, $j$ et $n$ tels que $P(X_n=i)\neq0$,
\[P_{(X_n=i)}(X_{n+1}=j)=t_{ij}.\]
Les nombres $t_{ij}$ sont les \emph{probabilités de transition}.
\end{definition}
\begin{definition}[Graphe probabiliste, matrice de transition]
La chaîne est représentée par un \emph{graphe probabiliste} : un graphe orienté pondéré dont les sommets sont les états, avec un arc de $i$ vers $j$, de poids $t_{ij}$, lorsque $t_{ij}>0$. La \emph{matrice de transition} est la matrice carrée $T=(t_{ij})$. Ses coefficients sont positifs et la somme de chaque ligne vaut $1$.
\end{definition}
\begin{definition}[Distribution]
La \emph{distribution} à l'instant $n$ est la matrice ligne $\pi_n=\begin{pmatrix}P(X_n=1)&\cdots&P(X_n=k)\end{pmatrix}$ ; $\pi_0$ est la distribution initiale.
\end{definition}
\begin{theorem}[Évolution de la distribution]
Pour tout entier naturel $n$ : $\pi_{n+1}=\pi_nT$ et $\pi_n=\pi_0T^n$.
\end{theorem}
\begin{proof}
Par la formule des probabilités totales, avec la partition $\left(X_n=1\right),\dots,\left(X_n=k\right)$ (on omet les événements de probabilité nulle, qui contribuent $0$) :
\[P(X_{n+1}=j)=\sum_{i=1}^{k}P(X_n=i)\,t_{ij},\]
qui est le coefficient $j$ de $\pi_nT$. La seconde formule s'obtient par récurrence.
\end{proof}
\begin{theorem}[Transition en $n$ étapes]
Si $P(X_0=i)=1$, alors pour tout état $j$, $P(X_n=j)$ est le coefficient de ligne $i$ et de colonne $j$ de $T^n$. Autrement dit, la probabilité de passer de l'état $i$ à l'état $j$ en $n$ transitions est $\left(T^n\right)_{ij}$.
\end{theorem}
\begin{proof}
La distribution initiale est la ligne $\pi_0$ qui a un $1$ en position $i$ et des $0$ ailleurs. Alors $\pi_n=\pi_0T^n$ est la ligne $i$ de $T^n$.
\end{proof}
\begin{definition}[Distribution invariante]
Une distribution $\pi$ (ligne de coefficients positifs de somme $1$) est \emph{invariante} si $\pi T=\pi$. Si $\pi_0=\pi$, alors $\pi_n=\pi$ pour tout $n$.
\end{definition}
\begin{theorem}[Cas de deux états]
Soit $T=\begin{pmatrix}1-a&a\\b&1-b\end{pmatrix}$ avec $0\le a\le1$, $0\le b\le1$ et $0<a+b<2$. La chaîne admet une unique distribution invariante $\pi=\begin{pmatrix}\frac b{a+b}&\frac a{a+b}\end{pmatrix}$, et quelle que soit la distribution initiale, $\pi_n$ converge vers $\pi$.
\end{theorem}
\begin{proof}
Écrivons $\pi_n=\begin{pmatrix}x_n&1-x_n\end{pmatrix}$. Alors $x_{n+1}=(1-a)x_n+b(1-x_n)=(1-a-b)x_n+b$. Cette suite arithmético-géométrique a pour point fixe $\ell=\frac b{a+b}$, et $x_n-\ell=(1-a-b)^n(x_0-\ell)$. Comme $-1<1-a-b<1$, $(1-a-b)^n\to0$ et $x_n\to\ell$. Une distribution invariante vérifie $x=(1-a-b)x+b$, donc $x=\ell$ : elle est unique.
\end{proof}
\begin{method}[Trois états : distribution invariante]
On résout le système formé par $\pi T=\pi$ et la condition $a+b+c=1$ ; l'une des trois équations de $\pi T=\pi$ est conséquence des deux autres. La convergence éventuelle de $\pi_n$ s'observe à la calculatrice.
\end{method}
\begin{pitfall}
La distribution est une ligne et l'on multiplie à droite par $T$ : $\pi_{n+1}=\pi_nT$, et non $T\pi_n$. Ce sont les lignes de $T$, et non ses colonnes, qui ont pour somme $1$.
\end{pitfall}

\mxnewpage
\section{Exercices}
\begin{exercise}[Du graphe à la matrice]
Un graphe probabiliste a trois sommets $1$, $2$, $3$ et les arcs suivants : $1\to1$ ($\frac12$), $1\to2$ ($\frac14$), $1\to3$ ($\frac14$), $2\to1$ ($\frac12$), $2\to3$ ($\frac12$), $3\to1$ ($\frac14$), $3\to2$ ($\frac14$), $3\to3$ ($\frac12$).
\begin{enumerate}
\item Écrire la matrice de transition $T$ et vérifier la somme des lignes.
\item Partant de l'état $1$, calculer la distribution après deux transitions.
\end{enumerate}
\end{exercise}
\begin{solution}
\begin{enumerate}
\item Chaque ligne a pour somme $1$ : \[T=\begin{pmatrix}\frac12&\frac14&\frac14\\\frac12&0&\frac12\\\frac14&\frac14&\frac12\end{pmatrix}.\]
\item $\pi_1=\begin{pmatrix}\frac12&\frac14&\frac14\end{pmatrix}$, puis $\pi_2=\pi_1T=\begin{pmatrix}\frac14+\frac18+\frac1{16}&\frac18+0+\frac1{16}&\frac18+\frac18+\frac18\end{pmatrix}=\begin{pmatrix}\frac7{16}&\frac3{16}&\frac38\end{pmatrix}$.
\end{enumerate}
\end{solution}

\begin{exercise}[Deux opérateurs]
Chaque mois, $\pc{10}$ des clients de l'opérateur $A$ passent chez $B$ et $\pc{20}$ des clients de $B$ passent chez $A$ ; les autres restent. Au départ, chaque opérateur a la moitié des clients. On note $x_n$ la part de $A$ après $n$ mois.
\begin{enumerate}
\item Écrire la matrice de transition et déterminer la distribution invariante.
\item Démontrer que $x_{n+1}=\num{0,7}x_n+\num{0,2}$, puis que $x_n=\frac23-\frac16\times\num{0,7}^n$.
\item Déterminer le plus petit $n$ tel que $x_n>\num{0,66}$.
\end{enumerate}
\end{exercise}
\begin{solution}
\begin{enumerate}
\item $T=\begin{pmatrix}\num{0,9}&\num{0,1}\\\num{0,2}&\num{0,8}\end{pmatrix}$. Avec $a=\num{0,1}$ et $b=\num{0,2}$, le cours donne $\pi=\begin{pmatrix}\frac23&\frac13\end{pmatrix}$.
\item $x_{n+1}=\num{0,9}x_n+\num{0,2}(1-x_n)=\num{0,7}x_n+\num{0,2}$. Le point fixe est $\frac23$ et $x_n-\frac23=\num{0,7}^n\left(\frac12-\frac23\right)=-\frac16\times\num{0,7}^n$.
\item $x_n>\num{0,66}$ équivaut à $\frac16\times\num{0,7}^n<\frac23-\num{0,66}=\frac1{150}$, soit $\num{0,7}^n<\num{0,04}$, soit $n>\frac{\ln\num{0,04}}{\ln\num{0,7}}\simeq\num{9,02}$. Le plus petit entier est $n=10$ ($\num{0,7}^9\simeq\num{0,0404}$ et $\num{0,7}^{10}\simeq\num{0,0282}$).
\end{enumerate}
\end{solution}

\begin{exercise}[Une marche à trois états]
Une particule se déplace sur trois cases alignées $1$, $2$, $3$. De la case $2$, elle va en $1$ ou en $3$ avec la probabilité $\frac14$ chacune et reste en $2$ avec la probabilité $\frac12$. D'une case extrême, elle reste sur place ou va en $2$, avec la probabilité $\frac12$ chacune.
\begin{enumerate}
\item Écrire la matrice de transition $T$.
\item Vérifier que $\begin{pmatrix}\frac14&\frac12&\frac14\end{pmatrix}$ est une distribution invariante.
\item Démontrer que c'est la seule.
\end{enumerate}
\end{exercise}
\begin{solution}
\begin{enumerate}
\item Dans l'ordre des cases $1$, $2$, $3$ : \[T=\begin{pmatrix}\frac12&\frac12&0\\\frac14&\frac12&\frac14\\0&\frac12&\frac12\end{pmatrix}.\]
\item $\frac14\times\frac12+\frac12\times\frac14=\frac14$ ; $\frac14\times\frac12+\frac12\times\frac12+\frac14\times\frac12=\frac12$ ; $\frac12\times\frac14+\frac14\times\frac12=\frac14$.
\item Avec $\pi=\begin{pmatrix}a&b&c\end{pmatrix}$, la première équation donne $\frac a2+\frac b4=a$, soit $b=2a$ ; la troisième donne de même $b=2c$. Avec $a+b+c=1$ : $a+2a+a=1$, d'où $a=c=\frac14$ et $b=\frac12$.
\end{enumerate}
\end{solution}

\begin{exercise}[Deux bassins]
On modélise la quantité d'eau (en milliers de litres) de deux bassins reliés par $U_n=\vc{x_n}{y_n}$ avec
\[U_{n+1}=AU_n+B,\qquad A=\begin{pmatrix}\num{0,6}&\num{0,2}\\\num{0,1}&\num{0,7}\end{pmatrix},\qquad B=\vc{10}{5}.\]
\begin{enumerate}
\item Déterminer la colonne $C$ telle que $C=AC+B$.
\item Exprimer $U_n$ en fonction de $A^n$, $U_0$ et $C$.
\item On admet que les coefficients de $A^n$ tendent vers $0$. Déterminer la limite de $U_n$. Dépend-elle de $U_0$ ?
\end{enumerate}
\end{exercise}
\begin{solution}
\begin{enumerate}
\item $(I_2-A)C=B$ avec $I_2-A=\begin{pmatrix}\num{0,4}&-\num{0,2}\\-\num{0,1}&\num{0,3}\end{pmatrix}$, de déterminant $\num{0,12}-\num{0,02}=\num{0,1}$. Donc $C=\frac1{\num{0,1}}\begin{pmatrix}\num{0,3}&\num{0,2}\\\num{0,1}&\num{0,4}\end{pmatrix}\vc{10}{5}=\vc{40}{30}$. Vérification : $\num{0,6}\times40+\num{0,2}\times30+10=40$ et $\num{0,1}\times40+\num{0,7}\times30+5=30$.
\item $U_n=A^n(U_0-C)+C$.
\item $U_n\to C=\vc{40}{30}$, quelle que soit la valeur de $U_0$.
\end{enumerate}
\end{solution}

\begin{exercise}[Transitions en trois étapes]
Un utilisateur d'une application est, chaque jour, actif ($1$) ou inactif ($2$), avec la matrice de transition $T=\begin{pmatrix}\num{0,7}&\num{0,3}\\\num{0,5}&\num{0,5}\end{pmatrix}$.
\begin{enumerate}
\item Calculer $T^2$ et $T^3$.
\item Un utilisateur est actif le jour $0$. Quelle est la probabilité qu'il soit inactif le jour $3$ ?
\item Déterminer la distribution invariante.
\end{enumerate}
\end{exercise}
\begin{solution}
\begin{enumerate}
\item $T^2=\begin{pmatrix}\num{0,64}&\num{0,36}\\\num{0,6}&\num{0,4}\end{pmatrix}$ et $T^3=\begin{pmatrix}\num{0,628}&\num{0,372}\\\num{0,62}&\num{0,38}\end{pmatrix}$.
\item C'est le coefficient $(1,2)$ de $T^3$ : $\num{0,372}$.
\item $a=\num{0,3}$ et $b=\num{0,5}$ : $\pi=\begin{pmatrix}\frac{\num{0,5}}{\num{0,8}}&\frac{\num{0,3}}{\num{0,8}}\end{pmatrix}=\begin{pmatrix}\frac58&\frac38\end{pmatrix}$.
\end{enumerate}
\end{solution}

\begin{exercise}[Un système de bonus-malus]
Un assureur classe chaque année ses clients en trois niveaux : bonus ($1$), base ($2$), malus ($3$). La matrice de transition annuelle est
\[T=\begin{pmatrix}\num{0,8}&\num{0,2}&0\\\num{0,6}&\num{0,2}&\num{0,2}\\0&\num{0,6}&\num{0,4}\end{pmatrix}.\]
\begin{enumerate}
\item Dessiner le graphe probabiliste.
\item Déterminer la distribution invariante.
\item Un client est au niveau de base. Quelle est la probabilité qu'il soit au niveau bonus deux ans plus tard ?
\end{enumerate}
\end{exercise}
\begin{solution}
\begin{enumerate}
\item Arcs $1\to1$ ($\num{0,8}$), $1\to2$ ($\num{0,2}$), $2\to1$ ($\num{0,6}$), $2\to2$ ($\num{0,2}$), $2\to3$ ($\num{0,2}$), $3\to2$ ($\num{0,6}$), $3\to3$ ($\num{0,4}$).
\item Avec $\pi=\begin{pmatrix}a&b&c\end{pmatrix}$ : $\num{0,8}a+\num{0,6}b=a$ donne $b=\frac a3$ ; $\num{0,2}b+\num{0,4}c=c$ donne $c=\frac b3=\frac a9$. Puis $a+\frac a3+\frac a9=1$, soit $a=\frac9{13}$, $b=\frac3{13}$, $c=\frac1{13}$.
\item Coefficient $(2,1)$ de $T^2$ : $\num{0,6}\times\num{0,8}+\num{0,2}\times\num{0,6}+\num{0,2}\times0=\num{0,6}$.
\end{enumerate}
\end{solution}

\begin{exercise}[Puissances de la matrice de transition à deux états]
Soit $T=\begin{pmatrix}1-a&a\\b&1-b\end{pmatrix}$ avec $a+b\neq0$. On pose $\lambda=1-a-b$, $Q=\frac1{a+b}\begin{pmatrix}b&a\\b&a\end{pmatrix}$ et $R=\frac1{a+b}\begin{pmatrix}a&-a\\-b&b\end{pmatrix}$.
\begin{enumerate}
\item Vérifier que $Q+\lambda R=T$, $Q^2=Q$, $R^2=R$ et $QR=RQ=0$.
\item En déduire par récurrence que $T^n=Q+\lambda^nR$ pour tout $n\ge1$.
\item On suppose $0<a+b<2$. Déterminer la limite de $T^n$ et interpréter ses lignes.
\end{enumerate}
\end{exercise}
\begin{solution}
\begin{enumerate}
\item Le coefficient $(1,1)$ de $Q+\lambda R$ vaut $\frac{b+(1-a-b)a}{a+b}=\frac{(a+b)-a(a+b)}{a+b}=1-a$ ; les autres se vérifient de même. Pour $Q^2$ : le coefficient $(1,1)$ vaut $\frac{b^2+ab}{(a+b)^2}=\frac b{a+b}$, et ainsi de suite. $QR$ a pour coefficient $(1,1)$ $\frac{ba-ab}{(a+b)^2}=0$ ; tous les coefficients de $QR$ et $RQ$ sont nuls.
\item Vrai pour $n=1$. Si $T^n=Q+\lambda^nR$, alors $T^{n+1}=(Q+\lambda^nR)(Q+\lambda R)=Q^2+\lambda QR+\lambda^nRQ+\lambda^{n+1}R^2=Q+\lambda^{n+1}R$.
\item $\abs\lambda<1$, donc $\lambda^n\to0$ et $T^n\to Q$. Les deux lignes de $Q$ sont égales à la distribution invariante : à long terme, l'état de départ est « oublié ».
\end{enumerate}
\end{solution}

\mxnewpage
\section{Problème}
\begin{problem}[Le modèle des urnes d'Ehrenfest]
Deux urnes $U$ et $V$ contiennent au total deux boules. À chaque étape, on choisit l'une des deux boules au hasard (de façon équiprobable) et on la change d'urne. On note $X_n$ le nombre de boules dans $U$ après $n$ étapes ; les états sont $0$, $1$, $2$. Au départ, $U$ contient les deux boules.
\begin{enumerate}
\item Justifier que $(X_n)$ est une chaîne de Markov, dessiner son graphe probabiliste et écrire sa matrice de transition $T$ (états dans l'ordre $0$, $1$, $2$).
\item Déterminer l'unique distribution invariante $\pi$. Comparer avec la loi du nombre de boules dans $U$ si chaque boule est placée au hasard dans l'une des deux urnes, indépendamment.
\item Calculer $\pi_1$, $\pi_2$, $\pi_3$, $\pi_4$. La suite $(\pi_n)$ converge-t-elle ?
\item Démontrer que pour tout $k\ge1$, $\pi_{2k}=\begin{pmatrix}\frac12&0&\frac12\end{pmatrix}$ et $\pi_{2k+1}=\begin{pmatrix}0&1&0\end{pmatrix}$, puis que $\frac12\left(\pi_{2k}+\pi_{2k+1}\right)=\pi$.
\item Expliquer, sans calcul, pourquoi $X_n$ a la même parité que $n$.
\end{enumerate}
\end{problem}
\begin{solution}[Problème]
\begin{enumerate}
\item L'étape suivante ne dépend que du nombre actuel de boules dans $U$. De l'état $0$ ou $2$, on passe à coup sûr en $1$ ; de l'état $1$, on passe en $0$ ou en $2$ avec la probabilité $\frac12$ chacun. \[T=\begin{pmatrix}0&1&0\\\frac12&0&\frac12\\0&1&0\end{pmatrix}.\]
\item $\pi T=\pi$ avec $\pi=\begin{pmatrix}a&b&c\end{pmatrix}$ donne $\frac b2=a$, $a+c=b$, $\frac b2=c$. Avec $a+b+c=1$ : $\pi=\begin{pmatrix}\frac14&\frac12&\frac14\end{pmatrix}$. C'est la loi binomiale de paramètres $2$ et $\frac12$.
\item $\pi_0=\begin{pmatrix}0&0&1\end{pmatrix}$, $\pi_1=\begin{pmatrix}0&1&0\end{pmatrix}$, $\pi_2=\begin{pmatrix}\frac12&0&\frac12\end{pmatrix}$, $\pi_3=\begin{pmatrix}0&1&0\end{pmatrix}$, $\pi_4=\begin{pmatrix}\frac12&0&\frac12\end{pmatrix}$. La suite alterne entre deux valeurs : elle ne converge pas.
\item $\begin{pmatrix}\frac12&0&\frac12\end{pmatrix}T=\begin{pmatrix}0&1&0\end{pmatrix}$ et $\begin{pmatrix}0&1&0\end{pmatrix}T=\begin{pmatrix}\frac12&0&\frac12\end{pmatrix}$ ; une récurrence sur $k$ conclut. La demi-somme vaut $\begin{pmatrix}\frac14&\frac12&\frac14\end{pmatrix}=\pi$.
\item Chaque étape modifie de $1$ le nombre de boules dans $U$, donc change sa parité. Au départ $X_0=2$ est pair.
\end{enumerate}
\end{solution}
\begin{teachernote}[Statut du problème]
Le modèle d'Ehrenfest est un « problème possible » du programme. Il montre qu'une distribution invariante peut exister sans que $\pi_n$ converge : la condition $0<a+b<2$ du cas à deux états a un analogue pour trois états, hors programme (chaîne apériodique). Prolongement : simuler la chaîne en Python avec $N$ boules et observer les fréquences des états sur une longue durée.
\end{teachernote}
